%% LyX 2.4.4 created this file.  For more info, see https://www.lyx.org/.
%% Do not edit unless you really know what you are doing.
\documentclass[british]{article}
\usepackage[T1]{fontenc}
\usepackage[utf8]{inputenc}
\usepackage{babel}
\usepackage{amsmath}

\begin{document}
\title{Inequality and Social Contests}
\maketitle

\section{Introduction }


\noindent

The relationship between inequality and status motivations remains theoretically unsettled. While status goods are often dismissed as socially wasteful (Veblen, 1899; Frank, 1985), empirical evidence aksi suggests that relative status improves subjective well-being by preserving a sense of dignity and position within one’s reference group (Clark, Frijters \& Shields, 2008). Whether such goods merely serve symbolic functions or actively shape economic behaviour has traditionally been treated as a sociological concern. However, status goods are not a mere symbolic reflection of socio-economic positions; they also constitute the interactions that build and maintain a social hierarchy. 

As hierarchy itself is subject to competing interpretations in both sociology and economics, the link between inequality and status motivation remains contested. Even though hierarchies appear universal, their acceptance and justifications vary significantly. Some have explained hierarchy through unequal access to material resources and structural advantages, while others emphasise upon reproduction of inequality through social networks, cultural knowledge, and norms. The current paper builds upon Weber’s (1922) notion of status groups and Goode’s (1978) commentary on consumption as a mechanism of social mobility -- situating status competitions within a closed economy of firms and labour, where economic interaction depends on adaptation of skills while also being constrained by slower-moving processes of social acclimatisation. 

In the model of fixed positions we consider, the individuals can readily invest in market-relevant adaptability -- such as skills, behavioural conformity, and professional signalling—but long-term social categorisation based on language, nationality, religion, or political identity evolves far more slowly and continues to shape trust, loyalty, and perceived legitimacy within organisations (Tilly, 1998; Bourdieu, 1986). Inequality may thus encourage positional competition and yet pose indirect costs through lack of potential growth in the closed economy. The paper uses the dynamic model a closed economy of labour and firm to explain both i) how social capital (adaptability) expenditure could change with rise in inequality and ii) how inequality may be a constraint in long-term growth. 

\section{Model}
\noindent

At the core of the model is a dichotomy between two forms of social positioning: adaptability and long-term categorisation. Adaptability refers to short-term, market-driven traits such as skills, linguistic competence, personality fit, and professional behaviour. These traits are shaped by transactional interactions and can be improved through social capital investments influenced by economic incentives. Long-term categorisation, by contrast, reflects slow-moving social identities such as nationality, race, religion, gender, and political affiliation. These are shaped through non-transactional, contractual, and trust-based processes of social cohesion and assimilation, and they evolve far more slowly than adaptability.

This dichotomy reflects a key dynamic in social life: individuals can readily adapt their behaviour to fit social and professional environments (e.g. the “chameleon effect”), but they also must remain constrained by social identities that shape how they interact in a social environment. While firms and institutions rely on adaptability to generate revenue, they also operate within broader social environments where long-term categorisation influences trust, loyalty, and stability.

The model integrates this dichotomy in Spence's labour-queue theory, in which firms rank workers based on perceived earning potential. While promotion may depend on both adaptability and categorisation, the firm's revenue is generated only through adaptability investments or expenditure. Workers therefore invest in both dimensions of social capital: adaptability to enhance short-term productivity and categorisation to secure long-term social positioning -- whereas firms only account for adaptability.

The model is formalised as a three-player game involving a firm and two workers competing in a labour queue (see Section \ref{sec:closed_economy}). The firm promotes the worker with higher perceived earning potential, while demoting or retaining the lower-ranked worker. This three-game formulation is a stylised view of a more general $N$-player economy where workers compete for positions with both social capital investments and social lottery (see Appendix for details). 

As a preview of the dynamics that follow, consider that when adaptability is irrelevant, promotion is based purely on slow-moving categorisation, and workers must wait their turn in a strict social hierarchy. When only adaptability matters, both firms and workers invest heavily in short-term skills. Lastly, when neither adaptability nor categorisation matter, all uncertainty as well as inequality disappears so that earnings are distributed equally without any labour sorting. 

\subsection{Closed Economy}\label{sec:closed_economy}

\noindent

The model assumes a closed economy in which all income circulates between workers and firm. Drawing on a Schumpeterian view, economic growth arises in the model from continual introductions of new or higher-quality goods. Workers supply labour and receive wages, and these wages constitute the basis of domestic demand. Firms use the revenue generated from this demand - either through investment financed by household savings - to fund innovation. Successful innovation then produces new varieties or improves existing products, generating the next round of income. This circular flow of wages, consumption, and innovation creates an endogenous engine of growth: consumption completes one innovation cycle by returning income to firms, while investment initiates the next by financing the search for new products. In long-run equilibrium, firms are assumed to have already chosen an optimal allocation of resources that sustains a steady stream of innovation and revenues.

The model for endogenous growth economy uses a notion of a socio-economic class by accommodating two mobility mechanisms - one through upskilling and another through a social lottery. While the former reflects the systematic skill-based organisation of labour market, the latter represents connections-driven growth of the employee. It is assumed that the average growth of a participant of the economy involves hedging one's bets in both the mechanisms. 

The essential feature of the model for worker's dual purposes of status rewards through social competitions and economic success is that the rewards through social networks are outside the firm's direct influence. In other words, not all factors that influence labour sorting criteria are under firm's control. Costs such as job search costs for workers, or influence of social networks are social constraints that the firm is obliged to work around. Since the firm benefits from higher wage income through innovation and revenue , the costs of upskilling do flow back to the firm - unlike the costs of building social connections (encompassing  slow-categorisation) that do not flow back to the firm. 




\subsection{Labour and Firm}\label{sec:firmbehaviour}

\noindent

In a labour queuing model, the income distribution would be set by earning potential that would eventually serve the maximisation of revenue for the firm - thus maintaining the endogenous growth with innovation and maintenance of assets. To explain the relationship between inequality and status motivations, we take a stylised view of the dynamics from the mobility scheme described in Section \ref{sec:nplayereconomy}. More specifically, we consider a scenario where both workers choose a certain mixing criterion $\lambda$ to distribute income between fixed lottery $L$ and pay costs towards upskilling imposed by the firm. The lottery cost $L$ is an exogenous cost of social connections that is not managed by the firm. The firm thus adjusts either or both of the pay-distribution ($\alpha$) or the cost of training $b$ to maximise its profits.

As in the model first proposed by Spence, the firm pays more to high-risk position. This is a social constraint necessitated by the need to value the firm's own skills-training that invests in. 


Formally, the following variables are set by workers and firm in the dynamic model

\begin{enumerate}
\item  $\alpha$ - income-sharing - The firm choses an arbitrary $0 < \alpha < \frac{1}{2} $  so that the top-performer is given $1-\alpha$ as wage while the bottom-of-the-queue performer is given $\alpha$ as wage
\item $\lambda$ - portion of wage spent on adaptability by worker 1
\item $\eta$ -  portion of wage spent on adaptability by worker 2
\end{enumerate}

The growth of assets in the economy is subject to asset-maintenance costs and a constant depreciation. The money towards adaptability that workers invest in flows back to firm but the money towards societal competitions does not flow immediately into the firm's asset-maintenance. All efficiency loss in collection of investments from workers are considered as part of asset-maintaining costs in the model. The depreciation of assets means that firm must relying on investments from the worker while the asset-maintenance costs imply that income addition doesn't grow assets linearly. In a steady-state equilibrium, a firm's income stream balances the decline in assets and costs incurred by maintenance of assets. The firm chooses $\alpha$ so that incoming revenue balances cost of maintenance and the depreciation. 

Formally, at any point in time $t$, the assets accumulated $A_t$ would decline with a depreciation $\delta$ over time while maintenance costs also apply. The firm distributes a fixed proportion $\kappa$ of  $A_t$ as wage - while the rest of $A_t(1-\kappa)$ contributes to existing assets. From $R_t\equiv \kappa A_t $, the provides $\alpha R_t$ to lower rank worker and $(1-\alpha)R_t$ to higher rank worker.

The low-rank worker chooses $\lambda_t$ so that $\alpha R_t \lambda_t$ is spent in adaptability and the higher-rank worker chooses $\eta_t$ so that $(1-\alpha)R_t \eta_t$ is spent on adaptability. The total contribution $\psi_t = R_t (\lambda_t \alpha+ (1-\alpha)\eta_t)$ flows back to the firm. The costs from distribution of $\kappa A_t$ i.e. the collection of $\psi_t$ are adjusted bearing a (power) cost-function  $K(\alpha)$ and maintaining assets $(1-\kappa)A_t$ bearing a cost-function $C(A_t)$.

The firm must thus choose $\alpha$ to maximise $A_{t+1}$ for given $A_t$ at every cycle where the dynamics is provided by $A_{t+1}=(1-\kappa-\delta) A_t + \kappa A_t (\lambda_t \alpha + (1-\alpha)\eta_t) - K(\alpha) - C(A_t)$ .

The optimisation for the low-performing worker who spends $\alpha R_t \lambda_t$ on adaptability and $\alpha R_t (1-\lambda_t)$ on categorisation. If the contest advantage is a fixed long-term subjective benefit $\Lambda$ for either workers, then  low-performing worker optimises $u(\lambda_t) \equiv \upsilon(\alpha R_t \lambda_t) +\Lambda \frac{\alpha R_t(1-\lambda_t)}{\alpha R_t(1-\lambda_t) + (1-\alpha) R_t(1-\eta_t)} $ whereas the high-performing worker optimises $U(\eta_t) \equiv \upsilon((1-\alpha) R_t \eta_t)  + \Lambda \frac{(1-\alpha) R_t(1-\eta_t)}{\alpha R_t(1-\lambda_t)+(1-\alpha) R_t(1-\eta_t)}$. The solutions for both can be found explicitly for $\upsilon(x)\equiv x$ and implicitly (numerical solution) using $\upsilon(x)\equiv \sqrt{x}$.








\section{Conclusions}
A higher lottery might reduce the sum of costs from all workers contributed to the firm. By changing stakes ($\alpha$) the firm may choose to decrease contributions or change the costs parameters ($\gamma$)  to raise collection. The 


\section{Appendix}


\subsection{The closed $N$-player economy } \label{sec:nplayereconomy}

\noindent


The stylised model for workers relies on the dynamics implied by an $N$-player economy where both social lottery and training costs suggest a separation of workers based on risks in the closed economy. In the microeconomic formulation of the endoegnous growth economy, the stratified labour is represented as fixed ladder of jobs occupied by workers at any given point in time. Each position in the ladder is associated with a distinct income level that is eventually set by the firm. Every worker holds exactly one job and the number of jobs does not change over time in the economy. While we assume no long-term unemployment, the workers in the economy face temporary job loss \footnote{The technological changes or internal reorganization that lead to such losses circusmscribe the creative destruction in endogenous growth economy.}. Upon the job-loss, the worker faces the choice to re-enter the job-market with costs of up-skilling and play a lottery to swap the current position with someone else in the job-market.

In other words, a displaced worker can return to the previous job by investing in the skills or certifications required to perform it. Since higher-ranked jobs with higher-income require greater skill, the upward move requires substantial retraining costs. The retraining costs are forgone by the worker in accepting lower-ranked jobs. Mobility thus occurs through three mechanisms i) voluntary movement down the ladder ii) costly investment to move upward iii) or a lottery for positions among other workers.

The model captures these possibilities with the three choices. The worker may  (i) swap her job with a lower-ranked worker who is willing to acquire the skills necessary to move up (therefore paying the price to move up) (ii) pay the skill cost required to claim a higher-ranked position from someone willing to step down, or (iii) pay a fixed cost that provides mobility with some risk (social lottery). Such options reflect different mechanisms of job mobility - training, competition, and luck.

The set up of the model also entails a variation of risk across incomes. More specifically, the higher-income positions have greater training requirements (factoring in a longer search time and other costs arising out of short-term unemployment) and thus a higher acquisition cost. Since workers in high-skill, high-income positions face greater absolute losses when displaced, the risk of income loss in absolute terms increases with higher income-levels. With all workers having the same preferences, the differences in behaviour across income levels are driven by well-known pattern of risk-aversion for small-gains and higher risk-taking for higher losses. This is captured by the model with an S-shaped Friedman-Savage utility for the worker.



Formally, the set up the model includes the following conditions. 
\begin{enumerate}

\item The choice for the worker is to either pay cost-of-upskilling or use the social lottery to move up the economic ladder. The lottery is of a fixed cost $L$ so that another worker $k$ out of N workers is randomly selected to swap her job with worker $i$. Worker $k$ does not pay any additional cost other than the fixed lottery cost $L$ for such a swap.

\item If not playing the lottery, a worker with income $m_i$  may take up a job that pays $m_{k}$ so long as the worker pays a transition-cost $C(m_i,m_{k})$ -- a power-cost-function defined as follows 

\[
C(m_i,m_k) =
\begin{cases}
(m_k - m_i)^b, & \text{for } m_i > m_k \\
0,            & \text{for } m_i < m_k
\end{cases}
\] Above ensures that $\forall {k\le i}$,  we have $C(m_i,m_{k})=0$ i.e. no cost is paid for moving from a position of high-skill to low (or the same) skill. Higher $b$ corresponds to higher costs (charged by the firm - as discussed in Section \ref{sec:firmbehaviour}).

\item The income-levels are specified with a single-parameter $a$ in the pareto income list $ x_0  (1 - i)^{-1 / a}$ for $1 \le i \le N$.  A higher $a$ corresponds to higher-stakes through income levels.

\end{enumerate}

The more a worker upskills oneself, the harder it becomes for the worker to move up (as the swap pool gets limited). A higher-income employee would also have less likelihood of being replaced (less transferrable-skills due to higher costs). The long-term transition-probability can thus be determined from the costs of moving up the ladder. 

It follows from the dynamics that in the long-term equilibrium, higher costs ($b$) discourage mobility and higher stakes ($a$) encourage bets.  The S-shaped utility means that in the long-term equilbrium, the higher-rank workers would prefer the downgrade while the low-rank workers may prefer the lottery. 



\end{document}
